\documentclass[10pt]{beamer}
\setbeamercolor{background canvas}{bg=white}
\usetheme[progressbar=frametitle,block=fill]{metropolis}

\usepackage{amsmath,amssymb,booktabs,array,graphicx,tabularx,multirow}
\usepackage{subcaption}
\usepackage{natbib}
\usepackage{tikz}
\usetikzlibrary{positioning,fit,calc,arrows.meta,decorations.pathreplacing,backgrounds}
\usepackage{xcolor,colortbl}

\definecolor{darkblue}{RGB}{12,44,92}
\definecolor{accent}{RGB}{153,0,0}
\definecolor{tcblue}{RGB}{31,119,180}
\definecolor{tcgreen}{RGB}{44,160,44}
\definecolor{tcorange}{RGB}{255,127,14}
\definecolor{softgray}{RGB}{240,242,245}

\setbeamercolor{frametitle}{fg=darkblue,bg=white}
\setbeamercolor{title}{fg=darkblue,bg=white}
\setbeamercolor{progress bar}{fg=darkblue,bg=softgray}
\setbeamercolor{alerted text}{fg=accent}
\setbeamerfont{title}{series=\bfseries}
\setbeamertemplate{navigation symbols}{}

\newcommand{\E}{\mathbb{E}}
\newcommand{\SR}{\mathrm{SR}}
\newcommand{\beaverarrow}{{\color{accent}$\hookrightarrow$\,}}
\newcommand{\myfontsize}{\fontsize{8}{9.5}\selectfont}
\newcommand{\backupbegin}{
   \newcounter{framenumberappendix}
   \setcounter{framenumberappendix}{\value{framenumber}}
}
\newcommand{\backupend}{
   \addtocounter{framenumberappendix}{-\value{framenumber}}
   \addtocounter{framenumber}{\value{framenumberappendix}}
}

\title{Liquidity-Aware Machine Learning\\for Stock Return Prediction}
\author{Daniele Bianchi$^{\dagger}$ \and Teng Jiao$^{\dagger}$}
\institute{Queen Mary University of London}
\date{[Conference Name], [Year]}

\begin{document}

\maketitle

%% ============================================================
%% MOTIVATION
%% ============================================================

\begin{frame}{ML Predictions Work, But Can You Trade Them?}
\begin{columns}[T]
\column{0.65\textwidth}
\textbf{The puzzle:}

\vspace{0.5em}
\begin{itemize}
  \item[\beaverarrow] ML achieves OOS $R^2 \approx 0.4\%$--$0.5\%$ for stock returns (GKX 2020)
  \item[\beaverarrow] But alpha concentrates in \textit{illiquid} stocks (Avramov et al.\ 2023)
  \item[\beaverarrow] After transaction costs, much of the alpha disappears
\end{itemize}

\vspace{0.6em}
\begin{block}{Our insight}
The problem starts at \textbf{training}, not portfolio formation. Standard ML weights all stocks equally --- but errors on untradeable stocks are costless.
\end{block}

\column{0.3\textwidth}
% Optional: small figure showing gross vs net SR
\end{columns}
\end{frame}


\begin{frame}{This Paper}
\textbf{Research question:} Does incorporating liquidity into ML \textit{training} improve implementable returns?

\vspace{1em}
\textbf{What we do:}
\begin{enumerate}
  \item Develop implementability-weighted training (cost-sensitive learning $\rightarrow$ asset pricing)
  \item Clean $2\times 2$ design separating training vs.\ portfolio effects
  \item Test across 3 ML models: XGBoost, Random Forest, Neural Network
\end{enumerate}

\vspace{1em}
\textbf{Key findings:}
\begin{itemize}
  \item[\beaverarrow] Training effect is \textbf{model-dependent}: helps XGBoost ($\Delta$SR $= +0.24$, $p = 0.033$), not RF
  \item[\beaverarrow] Mechanism: feature reallocation from illiquidity $\rightarrow$ tradable factors ($t = 4.55$)
  \item[\beaverarrow] Sequential boosting is vulnerable to implementability imbalance; bagging is not
\end{itemize}
\end{frame}


%% ============================================================
%% THE PROBLEM
%% ============================================================

\begin{frame}{The Implementability Imbalance}
\myfontsize
\begin{table}
\centering
\begin{tabular}{l ccccc}
\toprule
Quintile & Training & \$100M & \$500M & \$1B & \$5B \\
\midrule
Q1 (Illiquid) & 20\% & 11\% & 8\% & 7\% & 4\% \\
Q2 & 20\% & 19\% & 18\% & 17\% & 14\% \\
Q3 & 20\% & 22\% & 22\% & 22\% & 21\% \\
Q4 & 20\% & 24\% & 25\% & 26\% & 28\% \\
Q5 (Liquid) & 20\% & 25\% & 27\% & 29\% & 33\% \\
\bottomrule
\end{tabular}
\end{table}

\vspace{0.5em}
The model trains on a \textbf{uniform distribution} but must trade a \textbf{liquid-skewed} universe.

At \$500M: Q4--Q5 = 52\% of implementable stocks but only 40\% of training gradient.
\end{frame}


%% ============================================================
%% THE $2\times 2$ DESIGN
%% ============================================================

\begin{frame}{The $2\times 2$ Identification Strategy}
\begin{table}
\centering
\begin{tabular}{l cc}
\toprule
 & EW Portfolio & Liq-Weighted Portfolio \\
\midrule
Standard Training & \textbf{1A} (Baseline) & 1B \\
Weighted Training & 2A & \textbf{2B} (Combined) \\
\bottomrule
\end{tabular}
\end{table}

\vspace{1em}
\textbf{Effect decomposition:}
\begin{align*}
\text{Training Effect} &= \SR(2A) - \SR(1A) \\
\text{Portfolio Effect} &= \SR(1B) - \SR(1A) \\
\text{Total Effect} &= \SR(2B) - \SR(1A)
\end{align*}

If Training Effect $>$ Portfolio Effect $\Rightarrow$ \textit{how} you train matters more than \textit{how} you form portfolios.
\end{frame}


%% ============================================================
%% RESULTS
%% ============================================================

\begin{frame}{Main Results}
% \begin{figure}
% \centering
% \includegraphics[width=0.85\textwidth]{FiguresNew/effect_decomposition.png}
% \caption{Effect decomposition: Training vs.\ Portfolio effects by model}
% \end{figure}

\vspace{1em}
\textbf{XGBoost:} Training effect = $+0.242$, $p = 0.033$ (exceeds $\geq 0.20$ target)

\textbf{Random Forest:} Training effect = $-0.12$ (weighted training \textit{hurts})

\vspace{1em}
\begin{block}{Model dependence}
Sequential boosting (XGBoost) concentrates on illiquidity patterns $\rightarrow$ reallocation helps.\\
Bagging (RF) already diversifies $\rightarrow$ reallocation is redundant.
\end{block}
\end{frame}


\begin{frame}{Feature Reallocation (H2)}
% \begin{figure}
% \centering
% \includegraphics[width=0.85\textwidth]{FiguresNew/importance_comparison_xgboost.png}
% \end{figure}

\vspace{1em}
\textbf{XGBoost H2:} $t = 4.55$, $p < 0.001$ --- our strongest result.

Weighted training shifts SHAP importance: \\
\quad $\downarrow$ IdioVol, zero-trading days, MaxRet \\
\quad $\uparrow$ Mom12m, profitability, value
\end{frame}


%% ============================================================
%% CONCLUSION
%% ============================================================

\begin{frame}{Conclusion}
\begin{enumerate}
  \item Training-stage liquidity weighting can significantly improve implementable Sharpe ratios --- but \textbf{only for architectures vulnerable to implementability imbalance}
  \item The mechanism is feature reallocation: from illiquidity proxies to tradable factors
  \item Model architecture is a first-order consideration for cost-sensitive training
  \item Complements Jensen et al.\ (2024) [portfolio side] and Bianchi et al.\ (2026) [SDF side]
\end{enumerate}
\end{frame}


%% ============================================================
%% BACKUP
%% ============================================================

\backupbegin

\begin{frame}{Backup: Weighting Scheme}
The primary scheme is softmax on dollar volume rank:
$$w_i = \frac{\exp(\lambda \cdot \text{percentile}_i)}{\sum_j \exp(\lambda \cdot \text{percentile}_j)}, \quad \lambda = 2.0$$

Normalised to mean $= 1.0$ within each cross-section. All ML models accept \texttt{sample\_weight} natively.
\end{frame}

\begin{frame}{Backup: Transaction Cost Model}
Frazzini, Israel, and Moskowitz (2018):
$$\text{TC}_i = \frac{S_i}{2} + 0.1 \cdot \sigma_i \cdot \sqrt{\frac{Q_i}{\text{ADV}_i}}$$

\begin{itemize}
  \item Corwin-Schultz spreads $\times 0.30$ (institutional effective spread)
  \item Spread cap at 5\%
  \item AUM scenarios: \$100M, \$500M, \$1B, \$5B
\end{itemize}
\end{frame}

\backupend

\end{document}
